\documentclass{article}
\usepackage{amsmath}
\title{Sampled Total-Applied Joint Work Audit}
\author{UpstreamDrift}
\begin{document}
\maketitle
\section{Inputs and Assumptions}
For joint $j$, aligned SI samples at strictly increasing finite times $t_i$
provide applied torque $\tau_{ij}$ in N m and velocity $v_{ij}$ in rad/s.
Torque includes feedback and saturation. The joint ordering, coordinate signs
and frames must match. Compute $P_{ij}=\tau_{ij}v_{ij}$ in W and reconstruct
$P_j(t)$ linearly on each sample interval. This does not assume separately
linear torque and velocity, certify control hold semantics, or capture impulses.
\section{Work and Sign Separation}
Report, separately per joint,
\[
 W_j=\int P_j(t)\,dt,\quad
 W_j^+=\int\max(P_j(t),0)\,dt,\quad
 W_j^-=\int\min(P_j(t),0)\,dt.
\]
For endpoints $a,b$ over duration $h$, signed area is $h(a+b)/2$.
If both are nonnegative the positive area is the same; if both are
nonpositive it is zero. For opposite signs the exact positive triangle is
\[
 A^+=\frac{h}{2}\frac{\max(a,b)^2}{|a|+|b|}.
\]
Sum intervals and report $W_j^-=W_j-W_j^+$ and $\sum_j W_j^+$.
Thus positive work is accumulated before joint or temporal cancellation.
\section{Acceptance and Limits}
For $P=(2,-2)$ over two seconds, signed work is zero, positive work is
one joule and negative work is minus one joule. Unequal crossings and
nonuniform sample times are tested independently. Linear-power refinement
preserves each integral. Invalid dimensions, nonfinite inputs, nonincreasing
time and nonfinite derived power or work must raise errors.
The implementation is \texttt{scripts/audit\_optimization\_work.py}.
The README gives exact reproduction commands. Sampling and model uncertainty,
missing impulses, clinical safety and replay qualification remain outside this
audit; no optimization improvement or real-golfer acceptance is claimed.
\end{document}
